% RESEARCH_DRAFT_SHA256: 3fafedcbe4fb83f11ecd8a549b0b3ac9e6f2d97b807f6b36ba385f3ad1918323
% Layout and source-reference wrappers expanded for the web renderer.
\documentclass[11pt]{article}
\usepackage[margin=25mm]{geometry}
\usepackage{amsthm}
\usepackage{fontspec,unicode-math}
\setmainfont{Latin Modern Roman}
\setmathfont{Latin Modern Math}
\usepackage{amsmath,amssymb,mathtools}
\usepackage{xurl,hyperref}
\hypersetup{colorlinks=true,urlcolor=blue,linkcolor=blue}
\setcounter{secnumdepth}{0}
\setlength{\parindent}{0pt}
\setlength{\parskip}{5pt}
\setlength{\emergencystretch}{3em}
\newcommand{\sref}[2]{\hyperlink{source:#1:#2}{[S#2]}}
\newcommand{\cataloguescope}[1]{\paragraph{Source question.}#1}
\newcommand{\floor}[1]{\left\lfloor#1\right\rfloor}
\newtheorem{proposition}{Proposition}
\newtheorem{lemma}[proposition]{Lemma}
\begin{document}
\paragraph{Research manuscript.} OpenAI \texttt{gpt-6-astra}, reasoning \texttt{ultra}; 11 October 2026.
\paragraph{Preserved source body.} The complete catalogue mathematical text and sources follow with layout and source links expanded, including source qualifications and its historical preparation notes. New research begins at the section headed Research attempt.

\section{Quasilinear Finite-Field Polynomial Multiplication}\label{statement}

\subsection{Definitions and mathematical statement}

\paragraph{Shared notation.}Standard number systems are denoted by $\mathbb N,\mathbb Z,\mathbb Q,\mathbb R,\mathbb C$; index ranges, computational models and specialized notation follow the cited problem source and the definitions below.


$\mathbb F_p$ is the prime field, and multiplying two degree-at-most-$d$ inputs returns their exact product. $M_p(d)$ counts bit operations in the stated multitape Turing model. Uniform big-O means its implied constant is independent of both $p$ and $d$.
\paragraph{Statement recorded in the supplied source.}
For a prime $p$, let $M_p(d)$ denote the multitape-Turing-machine cost of multiplying polynomials of degree at most $d$ in $\mathbb{F}_p[x]$, and put $n=d\log p$. Prove that $M_p(d)=O(n\log n)$ uniformly in $p$ and $d$, including when $p=2$. \href{https://doi.org/10.4171/owr/2024/40}{[S2]}
\paragraph{Scope qualification.}
The source defines $n=d\log p$ as written. The degenerate zero-degree case is not silently redefined to use a different input-size expression.
\paragraph{Source status.}
The supplied catalogue labels this record open. Its source-status metadata was checked against the supplied snapshot on 2026-10-11. This import does not independently establish current open status or mathematical correctness. The catalogue literature assessment is dated 2026-08-22.
\subsection{Short English statement}
Multiply finite-field polynomials in the stated nearly linear bit complexity uniformly over field size and degree.
\subsection{Sources}
\begin{itemize}
\item[S1] Ulam AI and the UnsolvedMath Contributors, UnsolvedMath: A Curated Collection of Open Mathematics Problems, catalogue entry. Catalogue curation and metadata: CC BY 4.0. The supplied source question is retained; definitions and the English summary are editorial additions.
\url{https://huggingface.co/datasets/ulamai/UnsolvedMath}.
\item[S2] Explicit Methods in Number Theory (2025), Oberwolfach Reports, DOI 10.4171/owr/2024/40, pp. 2309--2368
\url{https://doi.org/10.4171/owr/2024/40}.
\item[S3] Harvey--van der Hoeven, Polynomial multiplication over finite fields
\url{https://doi.org/10.1145/3505584}.
\textit{Underlying Oberwolfach Reports material retains its source-specific licence, including CC BY-SA 4.0 where applicable; see THIRD\_PARTY\_NOTICES.md.}
\end{itemize}
\paragraph{Preparation and review.}
Original question attributed to the cited source and UnsolvedMath contributors. Codex (GPT-6; exact serving revision not exposed) assisted with definitions, English summaries, source inspection and formatting. No new human mathematical authorship or independent proof verification is asserted. Catalogue identity checked against the complete supplied registry and existing repository statements; known flagged repetitions and obvious identical conjectures excluded. Full mathematical-equivalence/literature review has not been performed. Statement preparation and final review followed contribution prompts C and D; automated checks verify structure and identity, not open status.
\label{end:statement}

\section{Research attempt}
\subsection{Literature, source context and scope}
The full original catalogue body and source qualifications are preserved above. The relevant primary workshop passages were read in the cached publisher report. Historical claims and subsequent attributed results are distinguished from the mathematical work below. A restricted example, standard lemma or known theorem is not counted as new progress on the entire original target. Missing assumptions and repeated or rhetorical extracts are not silently turned into different solved problems.

Harvey's workshop Conjecture 2 asks for a uniform bound in the bit-size parameter $n=d\log p$. The source already describes Kronecker substitution as easy in a large-characteristic regime. The author-hosted Harvey--van der Hoeven paper's Theorem 1.2 is conditional on a stated Linnik-constant hypothesis; its title alone is not an unconditional resolution. The following standard packing calculation is an attempted elementary reduction, not a new multiplication algorithm.
\subsection{Explicit packing and its small-characteristic cost}
Let $p$ be prime, $d\ge1$, and represent input coefficients of degree at most $d$ by integers in $\{0,\ldots,p-1\}$. Set
\[
 b=\left\lceil\log_2\big((d+1)(p-1)^2+1\big)\right\rceil,
 \quad B=2^b,\quad N=(d+1)b.
\]
\begin{proposition}
One multiplication of two integers of at most $N$ bits, followed by digit extraction and reduction modulo $p$, recovers the polynomial product. If integer multiplication costs $M(N)$ bit operations, the direct binary-remainder implementation costs
\[
 O\big(M(N)+db\log_2 p\big).
\]
For $p=2$ and $M(N)=O(N\log N)$, this gives $O(d(\log d)^2)$ for $d\ge2$, rather than the desired $O(d\log d)$.
\end{proposition}
\begin{proof}
Encode $a(x),c(x)$ as $a(B),c(B)$. Both are below $B^{d+1}$. The integer convolution coefficient $h_k=\sum_{i+j=k}a_ic_j$ satisfies $0\le h_k\le(d+1)(p-1)^2<B$. Therefore $a(B)c(B)=\sum_{k=0}^{2d}h_kB^k$ has exactly those coefficients as its base-$B$ digits, with no carries between digits. The digits occupy consecutive blocks of $b$ bits and can be read in $O(db)$ time.

To reduce each digit modulo $p$, scan its bits maintaining a remainder $0\le r<p$. Each next bit $e$ replaces $r$ by $(2r+e)\bmod p$, using one comparison and at most one subtraction on $O(\log p)$ bits. This proves the displayed upper bound. If $p=2$, reduction is simply the low bit of each block. Here $b=\lceil\log_2(d+2)\rceil$ and $N=\Theta(d\log d)$, yielding the final bound on substitution into $M$.
\end{proof}
\subsection{The exact failed step}
The tempting implication from optimal integer multiplication to optimal finite-field multiplication loses a logarithm because the packed input is $\Theta(d\log d)$ bits for fixed characteristic, not $\Theta(d)$. This is a failure of this reduction's cost analysis, not a lower bound against other algorithms. The source's sharper unconditional algorithms are not reproduced or improved here. The constant-degree boundary of the catalogue notation is preserved and is not used as a spurious asymptotic counterexample.
\section{Verification and remaining gaps}
Finite exact computations compare packed and direct convolution over several small prime fields, including the all-one binary input that reaches the central coefficient $d+1$. The missing step is a carryless or other uniform reduction avoiding the packing overhead without an unproved number-theoretic assumption.

\section{References and source attribution}
Primary workshop report: \url{https://doi.org/10.4171/owr/2024/40}. Relevant passages read: OWR 2024/40, printed pages 2339--2341, Conjecture 2 and conditional result.

\url{https://doi.org/10.1145/3505584}. Publisher DOI fetch returned HTTP403; used accessible author version instead.

\url{https://www.texmacs.org/joris/ffnlogn/ffnlogn.pdf}. Read introduction and Theorem1.2, explicitly conditional Linnik hypothesis; full algorithm proof not audited.

All retrievals were checked on 11 October 2026. Original question attribution remains in the preserved Sources section. External theorems are attributed and not claimed as new discoveries or fully verified merely by citation.


\paragraph{Integration review.} Codex (GPT-6; exact serving revision not exposed) reviewed the written arguments and expanded layout and source-reference wrappers. Repository metadata and the production web renderer were checked using Python 3.13 and JavaScript with MathJax 3.2.2. This model-assisted review is not independent human proof verification; no PDF compilation or formal proof checking is asserted.

\section{Completion Estimate}
\noindent\textbf{COMPLETION: 0\%}
No substantiated new progress toward the entire original target is claimed. The worked results and failed route are delimited above; a known result, illustrative model, or auxiliary obstruction is not counted as original-target completion.

\section{Final status and attribution}
\textbf{Attempt status: unresolved.} This describes this research attempt, not a new ruling on the source catalogue's status or on external literature claims.
The mathematics and this self-audit were generated by OpenAI \texttt{gpt-6-astra}, reasoning \texttt{ultra}, on 11 October 2026. Codex (GPT-6; exact serving revision not exposed) orchestrated the batch, read and reviewed the written mathematical arguments, independently reran the finite checks, and contributed formatting/repository review. No human mathematical authorship, independent human proof checking, or formal verification is claimed. The historical preparation paragraphs in the reproduced catalogue body concern that body only.

Contribution prompt A was applied to the complete definition and supplied source record, followed by the adversarial checks of prompt B. Local reproduction files are \texttt{batch026/checks.py}, \texttt{computations.json}, and the per-ID audit JSON. They were not public links at preparation time. Numerical checks are finite and do not prove unrestricted results. No LaTeX engine was found locally; compilation was not performed by this research agent.

\end{document}
